\chapter{Revisão da Literatura}
\label{chap_referencial}

Este capítulo está estruturado em duas partes: o referencial teórico, que apresenta os conceitos fundamentais para a compreensão do tema, e os trabalhos relacionados, que discutem as principais contribuições acadêmicas e aplicações recentes, articulando abordagens da literatura, comparações entre rotas tecnológicas e lacunas de pesquisa.

\section{Referencial teórico}
O referencial teórico desta dissertação está estruturado em três vertentes complementares: sustentabilidade e metas globais, tecnologias de combustíveis alternativos e adaptação do setor marítimo brasileiro. Em conjunto, essas dimensões fornecem a base conceitual necessária para compreender os desafios e as possibilidades da transição energética no transporte marítimo.

\subsection{Sustentabilidade e Metas Globais}
A discussão sobre a transformação do transporte marítimo está associada ao conceito de sustentabilidade e às agendas globais voltadas à mitigação das mudanças climáticas. A literatura mobiliza referenciais normativos e interpretativos que ajudam a compreender por que a transição energética do setor se tornou tema central no debate internacional.

O conceito de sustentabilidade proposto por \citet{elkington_cannibals_1998}, por meio do \emph{Triple Bottom Line}, constitui um referencial central. Ao articular os pilares econômico, ambiental e social, essa abordagem amplia a análise do desempenho do setor marítimo, indicando que decisões relacionadas a custos operacionais, eficiência energética e escolha de combustíveis não podem ser avaliadas isoladamente. Ao contrário, devem considerar, simultaneamente, os impactos ambientais associados às emissões atmosféricas e as implicações sociais decorrentes da atividade marítima. Esse enquadramento teórico legitima abordagens multicritério na avaliação de alternativas tecnológicas e energéticas para o setor.

As metas ambientais estabelecidas pela \gls{IMO}, particularmente aquelas relacionadas à redução do teor de enxofre nos combustíveis marítimos e à descarbonização progressiva até 2050, refletem esse alinhamento com agendas globais de sustentabilidade. Conforme discutido por \citet{van_global_2019}, as regulações da \gls{IMO} impõem limites rigorosos ao teor de enxofre, inicialmente em áreas de controle de emissões e, posteriormente, de forma global, o que desencadeia mudanças significativas nos processos de refino e na composição do \emph{mix} de combustíveis utilizados na navegação. Os autores destacam que essas medidas contribuem para reduções mensuráveis nas emissões de poluentes atmosféricos, ainda que gerem desafios adicionais relacionados a custos, investimentos e emissões indiretas de \gls{GEE} no processo de refino.

A literatura também enfatiza que tais esforços regulatórios devem ser compreendidos à luz de limites biofísicos mais amplos. \citet{mukherjee_perspective_2020}, ao discutirem o papel dos biocombustíveis e da \gls{CCS} no setor marítimo, ressaltam a necessidade de manter as atividades humanas dentro de limites ambientais compatíveis com a estabilidade do sistema terrestre. Essa perspectiva reforça a centralidade da redução das emissões de \gls{GEE} como elemento estruturante das estratégias de transição energética no transporte marítimo.

Nessa direção, o conceito de limites planetários, formulado por \citet{rockstrom_safe_2009}, figura entre os principais fundamentos científicos das metas globais de descarbonização. Essa abordagem estabelece fronteiras ecológicas críticas que não devem ser ultrapassadas para garantir a estabilidade e a resiliência do sistema terrestre, incluindo limites diretamente associados às mudanças climáticas e às emissões de \gls{GEE}. Ao reconhecer que o crescimento econômico e as atividades produtivas estão condicionados a restrições biofísicas, esse referencial dá base científica à imposição de metas ambientais rigorosas ao setor marítimo. Políticas de redução de emissões e de controle do uso de combustíveis fósseis no transporte marítimo podem, portanto, ser interpretadas não apenas como escolhas regulatórias, mas como respostas institucionais à necessidade de manter o sistema global de transportes dentro de limites ambientalmente seguros.

A discussão é complementada pela distinção entre sustentabilidade forte e fraca, proposta por \citet{daly_sustainable_2006}. Ao enfatizar que determinados recursos naturais são insubstituíveis, essa abordagem reforça as limitações de estratégias baseadas exclusivamente em compensações ou ganhos marginais de eficiência. No contexto do transporte marítimo, isso implica reconhecer que a transição energética não pode se restringir à redução incremental de emissões, mas deve incorporar escolhas que preservem o capital natural e reduzam a dependência estrutural de combustíveis fósseis.

Paralelamente aos fundamentos normativos da sustentabilidade, a literatura sobre transições sociotécnicas oferece aportes relevantes para compreender como essas transformações podem ocorrer na prática. A \gls{MLP}, desenvolvida por \citet{geels_technological_2002,geels_multi-level_2011}, interpreta as transições tecnológicas como processos evolutivos que envolvem a interação entre regimes estabelecidos, inovações emergentes e pressões externas, como mudanças regulatórias e demandas sociais. Sob essa ótica, a substituição do \emph{bunker fuel} por combustíveis alternativos representa a transformação de um regime sociotécnico consolidado, influenciado por regulações ambientais e dinâmicas de mercado.

\citet{schot_strategic_2008} destacam o papel das chamadas \emph{niches}, espaços protegidos que permitem o desenvolvimento e o amadurecimento de inovações sustentáveis. No setor marítimo, combustíveis alternativos, como os biocombustíveis líquidos, podem ser compreendidos como tecnologias de transição, cuja difusão depende de políticas específicas, experimentação controlada e processos de aprendizado que viabilizem reduções de custo e ganhos de escala. Estudos como o de \citet{international_council_on_clean_transportation_potential_2021} indicam que os biocombustíveis podem desempenhar papel relevante na redução das emissões do transporte marítimo, especialmente no curto e médio prazo, enquanto soluções como o hidrogênio ainda enfrentam limitações tecnológicas e de infraestrutura.

Por fim, a teoria da difusão de inovações, proposta por \citet{rogers_diffusion_2003}, contribui para explicar a adoção gradual dessas alternativas energéticas. A introdução progressiva de misturas de biocombustíveis, como B15 e B20, exemplifica um processo incremental, no qual ajustes técnicos, regulatórios e organizacionais permitem a assimilação das inovações pelo setor. Esse processo evidencia que a transição energética no transporte marítimo tende a ocorrer de forma gradual, condicionada por fatores técnicos, econômicos e institucionais.

\subsection{Instrumentos regulatórios europeus e seus efeitos sobre a descarbonização do transporte marítimo}

Embora esta dissertação tenha como foco o contexto brasileiro, a literatura indica que, no curto e médio prazo, vetores decisivos de mudança no transporte marítimo têm origem em iniciativas regulatórias regionais, com destaque para a União Europeia \citep{psaraftis_comparative_2021,eu_reg_2023_957,eu_reg_2023_1805}. Diferentemente das metas globais da \gls{IMO}, que definem objetivos agregados de longo prazo, os instrumentos europeus introduzem sinais econômicos e obrigações operacionais que afetam as decisões dos armadores, independentemente da bandeira do navio ou do local de abastecimento do combustível.

Nesse quadro, dois instrumentos merecem destaque: a inclusão do transporte marítimo no Sistema Europeu de Comércio de Emissões e a adoção do regulamento \emph{FuelEU Maritime}. Esses mecanismos atuam de forma complementar, ao precificar emissões absolutas de dióxido de carbono e, simultaneamente, impor limites à intensidade de \gls{GEE} dos combustíveis utilizados a bordo.

A inclusão do transporte marítimo no \gls{EUETS} representa mudança estrutural na governança climática do setor. Conforme discutido por \citet{psaraftis_comparative_2021}, mecanismos baseados em mercado têm a capacidade de internalizar externalidades ambientais de forma mais imediata do que metas voluntárias ou instrumentos puramente tecnológicos. A partir de 2024, navios que escalem portos da União Europeia passam a ser obrigados a adquirir permissões de emissão para cobrir parte de suas emissões de $CO_2$, com regime de implementação gradual que atinge 100\% das emissões elegíveis a partir de 2026. Para viagens entre portos europeus e portos de países terceiros, 50\% das emissões associadas ao trajeto são incluídas no sistema.

A literatura destaca que esse mecanismo cria custo direto associado ao uso de combustíveis fósseis convencionais, afetando a competitividade operacional das embarcações. \citet{psaraftis_comparative_2021} demonstram que a precificação do carbono no transporte marítimo tende a influenciar decisões de curto prazo, como velocidade, escolha de rotas e, sobretudo, composição do combustível, uma vez que essas decisões permitem reduzir emissões sem a necessidade imediata de investimentos elevados em novos sistemas de propulsão. Nessa perspectiva, o \gls{EUETS} não prescreve tecnologias específicas, mas incentiva economicamente a adoção de soluções capazes de reduzir emissões reportadas.

De forma complementar, o regulamento \emph{FuelEU Maritime}, em vigor a partir de 2025, estabelece limites progressivamente mais rigorosos para a intensidade de \gls{GEE} da energia utilizada a bordo dos navios que tocam portos da União Europeia. Diferentemente do \gls{EUETS}, que incide sobre emissões absolutas, o \emph{FuelEU Maritime} adota abordagem baseada em ciclo de vida, considerando as emissões \emph{well-to-wake}, ou seja, desde a produção do combustível até sua utilização no navio. Essa distinção é relevante, pois desloca o foco regulatório do controle de emissões locais para a avaliação integrada das cadeias energéticas.

Estudos como os de \citet{ce_delft_assessment_2020} indicam que combustíveis renováveis e de baixo carbono, especialmente combustíveis compatíveis com a infraestrutura existente (do inglês, \emph{drop-in fuels}) obtidos de resíduos, são favorecidos pela abordagem de intensidade de carbono ao longo do ciclo de vida promovida pelo regulamento \emph{FuelEU Maritime}, uma vez que podem ser adotados com a tecnologia existente e contribuem para o cumprimento de metas até 2030 antes que combustíveis de emissão zero, como hidrogênio e amônia, estejam amplamente disponíveis.

A combinação entre \gls{EUETS} e \emph{FuelEU Maritime} cria, portanto, duplo vetor de pressão regulatória: por um lado, penaliza economicamente o uso de combustíveis fósseis intensivos em carbono; por outro, estabelece padrões mínimos de desempenho climático para os combustíveis utilizados. Estudos como o de \citet{ce_delft_fueleu_2021} destacam que esse arranjo institucional tende a favorecer soluções de transição capazes de reduzir emissões de forma imediata, sem exigir transformações estruturais profundas na frota ou na infraestrutura portuária no curto prazo, ao alinhar os incentivos de mercado entre o preço do carbono no \gls{EUETS} e os requisitos de intensidade de carbono do \emph{FuelEU Maritime}.

Nesse contexto, os biocombustíveis líquidos surgem na literatura como alternativa intermediária entre o \emph{bunker fuel} e os combustíveis de emissão zero. Estudos como o de \citet{international_council_on_clean_transportation_potential_2021} mostram que combustíveis líquidos derivados de biomassa, como biodiesel de óleos residuais e diesel renovável, podem ser usados em motores marítimos existentes sem grandes modificações, por serem compatíveis com a infraestrutura existente, e apresentam redução de emissões ao longo do ciclo de vida quando comparados ao \emph{bunker fuel}. Estudos de avaliação de ciclo de vida também confirmam que certas rotas de biocombustíveis podem alcançar reduções significativas de \gls{GEE} em comparação com combustíveis fósseis, embora não sejam suficientes, por si só, para alcançar a neutralidade de carbono no longo prazo.

Do ponto de vista geográfico, um aspecto central desses instrumentos regulatórios europeus é que sua aplicação independe do local de abastecimento do combustível. As obrigações do \gls{EUETS} e do \emph{FuelEU Maritime} são acionadas pela escala do navio em portos da União Europeia, e não pela origem do combustível utilizado. Assim, combustíveis abastecidos em portos fora da UE, incluindo portos brasileiros, são considerados nos mecanismos de conformidade, desde que atendam aos critérios de rastreabilidade e elegibilidade estabelecidos pela regulação europeia.

Essa característica amplia o alcance dos instrumentos europeus para além de suas fronteiras territoriais, criando incentivos indiretos para a oferta de combustíveis marítimos de menor intensidade de carbono em portos extraeuropeus, em razão, por exemplo, do impacto da demanda europeia por biocombustíveis sobre mercados internacionais e cadeias de suprimento de combustíveis. Estudos como o de \citet{emsa_potential_2020} indicam que a demanda e a regulação de biocombustíveis na UE afetam fluxos comerciais globais e incentivos para fornecedores extra-regionais. No entanto, a efetividade desses mecanismos depende da robustez dos sistemas de certificação e da disponibilidade sustentável de matérias-primas, pois essas condições influenciam o desempenho ambiental e de ciclo de vida dos combustíveis \citep{cherubini_life_2011}.

\subsection{Tecnologias de Combustíveis Alternativos}
No plano técnico, a transição energética no transporte marítimo se materializa por meio da substituição progressiva do \emph{bunker fuel} por combustíveis de menor intensidade de carbono. Como esse processo é gradual e condicionado por restrições tecnológicas, econômicas e institucionais, a literatura combina interpretações sobre trajetórias tecnológicas, classificações institucionais das rotas energéticas e evidências sobre desempenho, viabilidade e sustentabilidade.

Nesse sentido, a contribuição clássica de \citet{schumpeter_business_2006} oferece primeiro referencial ao compreender a inovação como processo de \emph{destruição criativa}, no qual novos paradigmas tecnológicos substituem soluções estabelecidas, mesmo quando estas são amplamente difundidas e eficientes em seus próprios termos. Aplicado ao setor marítimo, esse raciocínio permite interpretar o abandono progressivo dos combustíveis fósseis tradicionais como resultado da convergência entre pressões regulatórias, ambientais e econômicas, e não apenas como decisão técnica isolada.

\citet{freeman_economics_1997} aprofundam essa leitura ao enfatizar o papel dos paradigmas tecnológicos e das trajetórias de inovação nos processos de transformação industrial. Sob essa perspectiva, os combustíveis alternativos não competem em condições homogêneas, mas se inserem em diferentes estágios de maturidade tecnológica. Os biocombustíveis, em particular, podem ser interpretados como inovações intermediárias, capazes de reduzir emissões e, simultaneamente, aproveitar parte da infraestrutura existente, criando condições para a futura difusão de soluções mais disruptivas, como o hidrogênio e a amônia. Essa abordagem reforça a ideia de que a transição energética no transporte marítimo ocorre em ondas sucessivas, nas quais tecnologias transitórias e de longo prazo coexistem.

Do ponto de vista técnico e institucional, diversas organizações têm proposto classificações que auxiliam na organização desse campo emergente. Relatórios do \gls{ICCT}, da \gls{IEA} e da \gls{DNV} convergem na identificação de cinco grandes categorias de combustíveis alternativos para o transporte marítimo: biocombustíveis, \emph{e-fuels} (combustíveis sintéticos produzidos a partir de hidrogênio e carbono capturado), gás natural liquefeito (\gls{LNG}), amônia e hidrogênio. Essa taxonomia é relevante porque permite comparar rotas tecnológicas distintas em termos de maturidade, potencial de redução de emissões, compatibilidade com motores existentes e necessidades de infraestrutura. Para a análise desenvolvida nesta dissertação, ela também ajuda a distinguir soluções mais próximas da adoção no curto prazo daquelas cuja difusão depende de transformações mais profundas em sistemas de propulsão, armazenamento e abastecimento.

No caso de combustíveis associados ao horizonte de longo prazo, o hidrogênio verde é apontado como alternativa promissora devido à sua capacidade de eliminar emissões diretas de carbono \citep{ustolin_extensive_2022}. Contudo, a literatura indica que barreiras tecnológicas, custos elevados e desafios de infraestrutura tornam sua adoção inviável no curto prazo, especialmente para grandes embarcações, como navios porta-contêineres. Esse cenário reforça a relevância de soluções intermediárias, capazes de reduzir emissões enquanto tecnologias mais disruptivas amadurecem.

Entre essas alternativas, os biocombustíveis têm recebido atenção especial devido ao seu potencial de aplicação no curto e médio prazo. Estudos indicam que o uso de biodiesel e outros biocombustíveis líquidos pode resultar em reduções significativas das emissões de dióxido de carbono em comparação ao \emph{bunker fuel}, além de apresentar desempenho energético compatível com a operação marítima \citep{kesieme_biofuel_2019,masum_comparing_2023,international_council_on_clean_transportation_potential_2021}. Ademais, há evidências de que motores marítimos podem operar com misturas de biodiesel de até B20 sem a necessidade de modificações estruturais relevantes, o que reduz barreiras técnicas à adoção dessas alternativas \citep{altarazi_effects_2022}.

No âmbito conceitual, o biobunker encontra respaldo na economia circular \citep{emaf_towards_2013}, ao enfatizar o uso de resíduos como matéria-prima e a redução da extração de recursos fósseis. A literatura aponta que o biodiesel produzido a partir de fontes residuais, como óleos de cozinha usados, pode reduzir significativamente emissões de \gls{GEE}, monóxido de carbono, hidrocarbonetos e material particulado, além de manter compatibilidade com motores diesel existentes \citep{lopresto_sustainable_2025}. Entretanto, desafios relacionados à estabilidade do combustível, à logística e à disponibilidade de \emph{feedstock} ainda representam barreiras à sua difusão em larga escala \citep{mohd_noor_biodiesel_2018}.

Apesar do potencial técnico, diversos estudos destacam limitações econômicas e institucionais para a adoção dos combustíveis alternativos. \citet{lindstad_reduction_2021} ressaltam os custos de produção e as incertezas de mercado como fatores críticos, especialmente quando comparados aos combustíveis fósseis convencionais. \citet{mukherjee_techno-economic_2023} argumentam que a efetividade das rotas alternativas de descarbonização no setor marítimo depende, em grande medida, de sua integração com tecnologias de \gls{CCS}, ampliando o impacto climático positivo dessas soluções.

Avanço importante na literatura recente é a combinação de análises técnico-econômicas (\gls{TEA}) com avaliações de ciclo de vida (\gls{ACV}), permitindo análise integrada de custos, emissões e impactos ambientais ao longo de toda a cadeia produtiva dos biocombustíveis. Estudos como os de \citet{tan_biofuel_2021} e \citet{li_techno-economic_2022} demonstram que essa abordagem é fundamental para evitar avaliações parciais, que considerem apenas as emissões na fase de uso e negligenciem impactos associados à produção, transporte e processamento do combustível.

Por fim, a viabilidade dos combustíveis alternativos deve ser analisada à luz da sustentabilidade do \emph{feedstock} e dos possíveis efeitos sistêmicos. \citet{zink_circular_2017} alertam para o risco do \emph{rebound effect}, no qual ganhos de eficiência podem ser parcialmente anulados pelo aumento da demanda energética. Nesse contexto, mecanismos de certificação e políticas públicas desempenham papel central na mitigação desses riscos, ao estabelecer critérios de rastreabilidade, intensidade de carbono e uso sustentável de matérias-primas. Iniciativas como a \emph{Renewable Energy Directive} II e III (RED II/RED III), na União Europeia, e o programa RenovaBio, no Brasil, ilustram a incorporação desses princípios no desenho regulatório.

Ao integrar os referenciais de inovação tecnológica, as classificações institucionais, as métricas de desempenho ambiental e as dimensões regulatórias, delineia-se um quadro consistente para compreender o papel dos combustíveis alternativos na substituição do \emph{bunker fuel}. Esse arcabouço evidencia que a transição energética no transporte marítimo é um processo multidimensional, no qual fatores técnicos, econômicos, ambientais e institucionais interagem de forma complexa.

\subsection{Adaptação do setor marítimo brasileiro}
A adaptação do setor marítimo brasileiro às mudanças regulatórias e tecnológicas discutidas nas seções anteriores depende da articulação entre três dimensões: (i) características estruturais do sistema portuário e da navegação, (ii) arcabouço institucional e regulatório que condiciona incentivos e padrões de conformidade e (iii) viabilidade logística e econômica das rotas de abastecimento (\emph{bunkering}) de combustíveis alternativos. Embora a agenda de descarbonização seja global, seus caminhos de implementação são territorializados, refletindo capacidades instaladas, infraestrutura portuária, disponibilidade energética e mecanismos de política pública.

\subsubsection{(i) Características estruturais do setor e desafios de transição}
A regulação internacional constitui vetor de pressão que se materializa em requisitos operacionais, como o limite global de enxofre em combustíveis marítimos \citep{imo_imo_2020}, estabelecido por emendas ao Anexo VI da \gls{MARPOL} e implementado globalmente a partir de 1º de janeiro de 2020. Essa medida reduz o limite de 3,5\% para 0,50\% m/m fora das Áreas de Controle de Emissão e mantém 0,10\% m/m nas \gls{ECA}, consolidando novo patamar regulatório para o mercado de combustíveis marítimos.

No caso brasileiro, a adaptação a essa dinâmica regulatória e tecnológica ocorre em ambiente no qual a infraestrutura portuária e a frota precisam se preparar não apenas para combustíveis em conformidade com o teor de enxofre, mas também para a introdução gradual de rotas alternativas, como hidrogênio, amônia e metanol renováveis, que exigem requisitos adicionais de armazenamento, manuseio, segurança e procedimentos de abastecimento. A análise dessas adaptações, do porto ao navio, é tratada de forma sistemática no estudo de \citet{wei_evaluating_2023}, que revisa aspectos críticos como compatibilidade de materiais, armazenamento em navios e portos e procedimentos de \emph{bunkering}, destacando que a discussão tecnológica precisa ser acompanhada por requisitos de infraestrutura e operação.

\subsubsection{(ii) Arcabouço institucional e regulatório no Brasil}
A viabilização de combustíveis alternativos no Brasil não é explicada apenas pela disponibilidade técnica, mas também por instrumentos de política pública capazes de criar sinais econômicos e critérios de sustentabilidade. Nesse sentido, a Política Nacional de Biocombustíveis (RenovaBio) tem relevância por estruturar metas de descarbonização associadas a instrumentos de certificação e crédito (\gls{CBIO}), ancorados em abordagem de ciclo de vida (\gls{ACV}), além de enfrentar desafios institucionais e regulatórios discutidos na literatura. \citet{grangeia_brazilian_2022} analisam o programa e apontam incertezas associadas à consolidação do arcabouço legal e regulatório, bem como debates sobre aspectos metodológicos e de elegibilidade, incluindo preocupações com rastreabilidade e deslocamento do uso da terra, o que é relevante quando se considera o uso de biocombustíveis como solução de transição também para o setor marítimo.

Complementarmente, \citet{pereira_renovacalc_2025} detalham o RenovaCalc como ferramenta computacional para cálculo de intensidade de carbono no âmbito do RenovaBio, reforçando a centralidade de métricas de ciclo de vida na governança climática dos biocombustíveis no Brasil. Essa dimensão é relevante para o setor marítimo porque combustíveis alternativos tendem a ser avaliados não apenas por emissões no uso, mas pela intensidade de carbono \emph{do poço ao tanque} e, em algumas abordagens, \emph{do berço ao túmulo}, o que condiciona sua legitimidade como opção de descarbonização.

\subsubsection{(iii) Capacidades e restrições para adoção de combustíveis alternativos nos portos}
No plano operacional, a transição energética marítima depende de portos capazes de oferecer infraestrutura e serviços de abastecimento compatíveis com novas rotas energéticas. O estudo de \citet{wei_evaluating_2023} sistematiza os requisitos para adaptação de navios e portos, enfatizando que as propriedades físico-químicas dos combustíveis e as necessidades de armazenamento e \emph{bunkering} alteram exigências de projeto e operação.

Em perspectiva aplicada ao Brasil, o relatório de \citet{icct_hubs_2025} avalia a prontidão de portos brasileiros para apoiar a produção, o abastecimento (\emph{bunkering}) e o uso de combustíveis marítimos renováveis derivados de hidrogênio, como amônia e metanol, no contexto de corredores verdes. O relatório identifica portos candidatos e propõe critérios de prontidão relacionados à infraestrutura, à conectividade e às condições para suportar cadeias de suprimento de combustíveis renováveis, reforçando que a transição envolve decisões de investimento e coordenação entre atores públicos e privados.

\subsubsection{(iv) Biocombustíveis como rota de transição e o papel da logística e do feedstock}
Entre as rotas tecnológicas discutidas no item 2.1.2, os biocombustíveis líquidos destacam-se como alternativas potencialmente mais próximas da adoção por aproveitarem parte da infraestrutura existente. Para o caso brasileiro, a viabilidade desse caminho depende de condições logísticas e de suprimento de biomassa e resíduos em escalas compatíveis com demandas portuárias. \citet{carvalho_biofuels_2021} avaliam o potencial e a logística de biocombustíveis marítimos (biobunkers) considerando dimensões espacial, técnico-econômica e logística, incluindo a relação entre disponibilidade de matérias-primas e proximidade de infraestrutura portuária. Esse enfoque contribui para deslocar a análise de uma avaliação abstrata de \emph{potencial energético} para avaliação territorializada, na qual custos, distâncias e infraestrutura condicionam a viabilidade de adoção.

A adaptação do setor marítimo brasileiro deve, assim, ser compreendida como um processo que articula pressão regulatória internacional, como a decorrente de \citep{imo_imo_2020}, capacidade portuária e logística para novos combustíveis e instrumentos domésticos capazes de gerar critérios de sustentabilidade e previsibilidade econômica, como os mecanismos de \gls{ACV} e certificação associados ao RenovaBio. Em linha com os fundamentos apresentados no item 2.1.1 e com as rotas tecnológicas discutidas no item 2.1.2, o caso brasileiro evidencia que a transição energética marítima não se reduz a uma escolha tecnológica: trata-se de transformação institucional e infraestrutural que depende de coordenação, mensuração robusta da intensidade de carbono e planejamento portuário orientado à implantação progressiva de alternativas.

\section{Trabalhos relacionados}

\subsection{Estratégia de busca e critérios de seleção dos estudos}
Para organizar e analisar de forma sistemática os estudos selecionados, adota-se o modelo \gls{PECA}, utilizado em revisões estruturadas quando o foco da revisão envolve exposição e comparação entre alternativas \citep{morgan_identifying_2018}. No contexto desta pesquisa, as dimensões do \gls{PECA} são adaptadas da seguinte forma: (a) \emph{População}, representando o setor de transporte marítimo, em escala global e regional; (b) \emph{Exposição}, correspondente às rotas e tecnologias de suprimento e uso de combustíveis no lugar do \emph{bunker fuel}, com ênfase em biocombustíveis líquidos e misturas de biodiesel; (c) \emph{Comparação}, envolvendo a comparação entre combustíveis fósseis, biocombustíveis e combustíveis emergentes, como hidrogênio, amônia e \emph{e-fuels}; e (d) \emph{Avaliação}, contemplando evidências técnicas, econômicas e ambientais, incluindo custos, desempenho operacional e redução de emissões de \gls{GEE}.

A busca bibliográfica é realizada por meio da plataforma \emph{Elicit}, ferramenta baseada em técnicas de recuperação de informação por similaridade semântica e aprendizado de máquina, que permite identificar estudos relevantes mesmo quando estes não utilizam exatamente os mesmos termos da pergunta de pesquisa. Diferentemente de buscas tradicionais baseadas exclusivamente em palavras-chave, o \emph{Elicit} possibilita a triagem de artigos a partir de sua proximidade conceitual com a questão investigada, além de oferecer funcionalidades de extração estruturada de informações metodológicas e resultados \citep{lau_comparison_2025,bernard_using_2025}.

A pergunta de pesquisa utilizada como base para a busca é formulada da seguinte forma: \emph{What evidence exists on the technical, economic, and environmental performance of biofuels (e.g., biobunker blends) in maritime transport compared to conventional fossil fuels or emerging zero-emission alternatives?}. Essa formulação permite capturar estudos que abordam não apenas o desempenho ambiental dos biocombustíveis, mas também suas implicações técnicas, econômicas e comparativas em relação a outras rotas de descarbonização.

A busca resulta inicialmente em 37 artigos científicos e relatórios técnicos, publicados entre 2010 e 2024, provenientes de bases consolidadas da literatura científica internacional, como Scopus, Web of Science e ScienceDirect, além de documentos institucionais de referência elaborados por organismos como o \gls{ICCT}, a \gls{IEA} e a própria \gls{IMO}. Esses estudos cobrem amplo espectro de abordagens metodológicas, contextos geográficos e tecnologias energéticas aplicadas ao setor marítimo.

Dada a inviabilidade de análise aprofundada de todos os estudos identificados, procede-se a etapa de triagem e priorização, orientada por critérios de relevância temática, metodológica e substantiva em relação aos objetivos desta dissertação. São priorizados trabalhos que atendem a pelo menos um dos seguintes critérios: (i) avaliação direta do uso de biocombustíveis líquidos no transporte marítimo; (ii) comparação entre biocombustíveis, combustíveis fósseis convencionais e combustíveis emergentes de emissão zero; (iii) emprego de \gls{ACV}, \gls{TEA} ou abordagens integradas \gls{TEA}--\gls{ACV}; (iv) incorporação de aspectos logísticos, territoriais ou de infraestrutura de abastecimento (\emph{bunkering}); e (v) relevância explícita para contextos de países emergentes ou para o caso brasileiro.

Com base nesses critérios, seleciona-se um conjunto central de aproximadamente vinte estudos considerados essenciais para sustentar a análise dos trabalhos relacionados. Os demais artigos identificados na busca inicial são utilizados de forma complementar, como apoio conceitual ou empírico pontual, sem constituir o núcleo da revisão.

\subsection{Contexto e escopo dos estudos analisados}
Os estudos selecionados para esta seção abrangem amplo espectro de contextos geográficos, tecnológicos e institucionais relacionados à transição energética no transporte marítimo, refletindo a natureza global do setor e, simultaneamente, as especificidades regionais que condicionam a adoção de combustíveis alternativos. De modo geral, a literatura analisada pode ser organizada em três grandes níveis de escopo: (i) análises globais e setoriais voltadas às metas internacionais de descarbonização, (ii) estudos comparativos de rotas tecnológicas e energéticas e (iii) investigações aplicadas a contextos regionais ou nacionais, com destaque para países em desenvolvimento, particularmente o caso brasileiro.

No plano global, parte significativa da literatura concentra-se em avaliar as implicações das metas climáticas estabelecidas pela \gls{IMO} e sua interação com diferentes rotas de descarbonização. Trabalhos de referência elaborados pelo \gls{ICCT}, como os de \citet{international_council_on_clean_transportation_potential_2021}, estruturam o debate ao analisar o potencial de redução de emissões associado a combustíveis alternativos líquidos, incluindo biocombustíveis, em comparação ao \emph{bunker fuel}. Esses estudos adotam perspectiva sistêmica, considerando tanto emissões diretas quanto indiretas ao longo do ciclo de vida, e enfatizam que os biocombustíveis líquidos representam uma das poucas opções tecnicamente viáveis para reduções significativas de emissões no curto e médio prazo, especialmente até o horizonte de 2030.

Complementarmente, \citet{kass_understanding_2018} e \citet{mukherjee_perspective_2020} ampliam esse enquadramento ao discutir os biocombustíveis no contexto de portfólio mais amplo de tecnologias de descarbonização, incluindo \emph{e-fuels}, hidrogênio e soluções baseadas em \gls{CCS}. Esses autores argumentam que a transição energética no transporte marítimo não deve ser interpretada como substituição imediata e completa de combustíveis fósseis por alternativas de emissão zero, mas como processo gradual, no qual tecnologias de transição desempenham papel estratégico ao reduzir emissões enquanto soluções de longo prazo amadurecem técnica e economicamente. Essa perspectiva dialoga diretamente com a abordagem de transições sociotécnicas e com a ideia de coexistência de rotas energéticas ao longo do tempo.

No que se refere ao escopo tecnológico, diversos estudos adotam abordagem comparativa entre diferentes combustíveis alternativos. A literatura indica que a avaliação da disponibilidade, dos custos e do potencial de mitigação de emissões dos combustíveis alternativos é sensível aos instrumentos de política climática considerados, como precificação de carbono e mandatos de uso. De forma semelhante, \citet{krantz_assessing_2023} comparam oito combustíveis alternativos utilizando diferentes metodologias de inventário de emissões, demonstrando que escolhas metodológicas, como a distinção entre emissões médias e marginais, influenciam a avaliação do impacto climático das rotas energéticas analisadas.

Outros estudos ampliam o escopo da análise ao integrar o transporte marítimo a debates mais amplos sobre a transição energética nos setores de transporte. \citet{arias_assessing_2024}, por exemplo, projetam o desenvolvimento de tecnologias emergentes tanto para a aviação quanto para o transporte marítimo, destacando possíveis sinergias e disputas por recursos energéticos, especialmente no que diz respeito à disponibilidade de hidrogênio e combustíveis sintéticos. Essa abordagem intersetorial reforça a relevância dos biocombustíveis líquidos, que podem se beneficiar de cadeias produtivas existentes e reduzir pressões competitivas sobre rotas tecnológicas ainda incipientes.

No âmbito institucional, a literatura aponta que a viabilização dos combustíveis alternativos depende de condicionantes que extrapolam o desempenho técnico e ambiental. \citet{prussi_potential_2021} destacam o papel de critérios de sustentabilidade, regras de elegibilidade e prioridades regulatórias que moldam a disponibilidade e o uso de biocombustíveis. \citet{harahap_renewable_2023}, por sua vez, incorporam instrumentos institucionais, como metas climáticas e precificação de carbono, como parâmetros de cenário na modelagem de cadeias de suprimento. Em perspectiva mais ampla, \citet{irena_advanced_2019} evidenciam que a previsibilidade e a estabilidade regulatória, especialmente no desenho de mandatos, critérios de elegibilidade e políticas de apoio a biocombustíveis avançados, são determinantes centrais das decisões de investimento e da difusão desses combustíveis nos setores de transporte, pois a incerteza normativa representa uma das maiores barreiras à expansão da indústria.

Quando se desloca o foco para contextos regionais e nacionais, observa-se lacuna importante na literatura, especialmente no que se refere a países emergentes. Embora existam numerosos estudos globais e europeus, poucos trabalhos incorporam explicitamente as características territoriais, logísticas e institucionais de economias como a brasileira. Nesse sentido, \citet{carvalho_biofuels_2021} representam contribuição central ao analisar a viabilidade logística e técnico-econômica de biocombustíveis marítimos a partir de abordagem espacial, avaliando a disponibilidade de biomassa e os fluxos de suprimento até potenciais pontos de produção e abastecimento portuário. O estudo evidencia que a viabilidade do biobunker depende da proximidade entre fontes de biomassa, infraestrutura portuária e rotas de navegação, deslocando a análise de plano abstrato para contexto territorializado.

No mesmo sentido, \citet{cesilla_de_souza_technical-economic_2024} avançam ao integrar análise técnico-econômica e \gls{ACV} na comparação de biocombustíveis navais derivados de diferentes \emph{feedstocks} brasileiros, como sebo bovino e macaúba. Ao aplicar abordagem integrada considerando parâmetros representativos de embarcação marítima, os autores demonstram o potencial de redução de emissões dessas rotas e evidenciam a sensibilidade dos resultados a fatores como custo da matéria-prima, eficiência de conversão e escala de produção. Esses achados reforçam que a viabilidade ambiental e econômica dos biocombustíveis marítimos depende do contexto de aplicação, ressaltando a importância de análises aplicadas e contextualizadas para a compreensão dos desafios associados à transição energética no setor marítimo.

De forma geral, o escopo dos estudos analisados revela literatura robusta no que diz respeito à identificação de rotas tecnológicas e à avaliação comparativa de combustíveis alternativos, mas ainda fragmentada quanto à integração entre desempenho técnico, custos, logística e contexto institucional, especialmente em países emergentes. Observa-se predominância de análises agregadas e cenarizadas, com menor atenção a dados operacionais reais, simulações por rota ou embarcação e mecanismos nacionais de governança climática.

\subsection{Abordagens metodológicas adotadas na literatura}
A literatura sobre combustíveis alternativos no transporte marítimo apresenta diversidade significativa de abordagens, refletindo a complexidade técnica, econômica e ambiental envolvida nos processos de descarbonização do setor. De modo geral, os estudos analisados podem ser organizados em quatro grandes vertentes metodológicas: (i) \gls{ACV}, (ii) \gls{TEA}, (iii) abordagens integradas \gls{TEA}--\gls{ACV} e (iv) modelagem operacional, logística e espacial. A \gls{ACV} é mobilizada sobretudo para estimar impactos ambientais ao longo da cadeia de produção e uso do combustível; a \gls{TEA}, para examinar custos, escalas e viabilidade econômica; as abordagens integradas, para articular essas duas dimensões; e a modelagem operacional, logística e espacial, para aproximar a análise das condições reais de abastecimento, infraestrutura e operação. Cada uma dessas vertentes oferece contribuições específicas, mas também apresenta limitações que influenciam a interpretação dos resultados e a formulação de políticas e estratégias de transição energética.

\subsection{Avaliação de Ciclo de Vida (ACV)}
A \gls{ACV} constitui uma das abordagens mais utilizadas na literatura para analisar o desempenho ambiental de combustíveis alternativos no transporte marítimo. Essa abordagem permite quantificar emissões e impactos ambientais ao longo de toda a cadeia de produção, transporte e uso do combustível, superando análises restritas à fase de combustão a bordo. Diversos estudos destacam a relevância da \gls{ACV} para evitar conclusões enviesadas que desconsiderem emissões indiretas associadas à produção de matérias-primas, processamento e logística.

Trabalhos de referência, como os de \citet{international_council_on_clean_transportation_potential_2021}, aplicam modelos consolidados de avaliação de ciclo de vida para comparar combustíveis fósseis, biocombustíveis líquidos e outras rotas alternativas, indicando que rotas específicas de biocombustíveis avançados, especialmente aquelas baseadas em resíduos lipídicos e biomassa lignocelulósica, podem apresentar reduções expressivas de emissões de \gls{GEE} quando avaliadas sob perspectiva de ciclo de vida abrangente. De forma semelhante, \citet{bilgili_comparative_2021} realizam avaliação comparativa em perspectiva de ciclo de vida, mostrando que biocombustíveis derivados de resíduos tendem a apresentar melhor desempenho ambiental em comparação a rotas baseadas em culturas dedicadas, especialmente quando considerados impactos relacionados ao uso da terra.

Outros estudos adotam inventários próprios ou bases de dados específicas para refinar as estimativas de impacto ambiental. \citet{masum_comparing_2023}, ao aplicar \gls{ACV} diretamente ao setor marítimo, desenvolvem inventários detalhados para diferentes rotas de biocombustíveis líquidos, evidenciando a sensibilidade dos resultados a características tecnológicas e de \emph{feedstock}. De forma complementar, \citet{jeswani_environmental_2020}, em revisão crítica sobre a sustentabilidade ambiental dos biocombustíveis, demonstram que resultados de \gls{ACV} podem variar significativamente em função de escolhas metodológicas, como limites do sistema, fatores de alocação e fontes de dados, reforçando a necessidade de transparência e padronização na aplicação da metodologia.
Cabe destacar que a literatura distingue duas abordagens principais na aplicação da Avaliação de Ciclo de Vida: a abordagem attributional e a consequential. A ACV attributional busca estimar os impactos médios associados a um sistema produtivo em determinado estado, sendo amplamente utilizada em análises comparativas de tecnologias sob condições estáticas. Por sua vez, a ACV consequential procura capturar os efeitos marginais decorrentes de mudanças na demanda, incorporando ajustes no sistema produtivo, como expansão de oferta, substituição de insumos e mudanças indiretas no uso da terra.

No contexto desta dissertação, adota-se a abordagem attributional, alinhada ao objetivo de comparar cenários sob condições operacionais constantes e com base em fatores médios de emissão extraídos da literatura. Ainda assim, a distinção é relevante, pois evidencia que resultados baseados em médias podem não refletir integralmente os impactos associados à expansão em larga escala de determinadas rotas de biocombustíveis, conforme discutido na literatura recente.

Um dos principais mecanismos associados a essas dinâmicas é a mudança indireta do uso da terra (\gls{ILUC}, do inglês \emph{Indirect Land Use Change}). O ILUC ocorre quando o aumento da demanda por matérias-primas agrícolas destinadas à produção de biocombustíveis desloca outras atividades, como a produção de alimentos, para novas áreas, podendo resultar na conversão de ecossistemas naturais, como florestas ou pastagens, em áreas agrícolas.

Esse processo pode gerar emissões adicionais de gases de efeito estufa associadas à mudança no estoque de carbono do solo e da vegetação, que não são capturadas em abordagens \emph{attributional}, mas podem ser parcialmente incorporadas em análises \emph{consequential}. Assim, embora determinadas rotas de biocombustíveis apresentem baixos fatores de emissão médios, seus impactos climáticos totais podem ser significativamente maiores quando considerados efeitos indiretos de uso da terra.

A literatura aponta que a magnitude do ILUC varia de acordo com o tipo de matéria-prima, a região de produção e as condições de mercado, sendo particularmente relevante para culturas associadas à expansão agrícola, como a soja. Esse aspecto reforça a necessidade de cautela na interpretação de resultados de \gls{ACV} baseados exclusivamente em valores médios, especialmente em análises que buscam avaliar a viabilidade de expansão em larga escala de biocombustíveis.

Apesar de sua robustez conceitual, a literatura também aponta limitações da avaliação de ciclo de vida quando aplicada de forma isolada à análise da descarbonização do transporte marítimo. Estudos recentes evidenciam que abordagens baseadas apenas em emissões médias podem ocultar efeitos marginais relevantes, especialmente em cenários de expansão ou substituição em larga escala de combustíveis. Nesse sentido, \citet{krantz_assessing_2023}, ao comparar oito combustíveis alternativos por meio de análises de ciclo de vida \emph{attributional} e \emph{consequential}, demonstram que, embora biocombustíveis apresentem baixos valores médios de emissões, suas emissões marginais podem ser mais elevadas em função da limitação de matérias-primas sustentáveis. Ademais, os autores ressaltam que a \gls{ACV}, por si só, não captura adequadamente dimensões econômicas, operacionais ou institucionais, o que restringe sua capacidade de orientar decisões de investimento ou estratégias de transição quando utilizada de forma isolada. Em conjunto, esses resultados indicam que transições baseadas exclusivamente em biocombustíveis líquidos podem não produzir os benefícios climáticos esperados, reforçando a importância de considerar efeitos sistêmicos e restrições de oferta na análise da descarbonização do transporte marítimo. Essa limitação ajuda a explicar por que parte da literatura avança da avaliação ambiental para abordagens orientadas à viabilidade econômica e, posteriormente, para modelos integrados.

\subsection{Análise Técnico-Econômica (TEA)}

A \gls{TEA} constitui outra vertente central da literatura e é utilizada para avaliar a viabilidade financeira e os custos associados às diferentes rotas de combustíveis alternativos. A \gls{TEA} busca estimar custos de produção, capital e operação, permitindo comparar alternativas sob a ótica econômica a partir de indicadores como custo total e custo por viagem.

Estudos como os de \citet{tan_biofuel_2021}, \citet{li_techno-economic_2022} e \citet{tan_economic_2019} aplicam \gls{TEA} para comparar diferentes rotas de produção de biocombustíveis, considerando parâmetros como eficiência de conversão, custos de matéria-prima, escala de produção e maturidade tecnológica. Esses trabalhos indicam que, embora os biocombustíveis líquidos apresentem custos superiores aos combustíveis fósseis em cenários de mercado não regulados, sua viabilidade econômica tende a melhorar quando considerados mecanismos de política climática, como precificação de carbono ou incentivos à produção.

No contexto marítimo, a \gls{TEA} também é utilizada para comparar combustíveis alternativos entre si, incluindo \emph{e-fuels}, hidrogênio e amônia. \citet{lindstad_reduction_2021} demonstram que combustíveis de emissão zero apresentam custos superiores no curto prazo, o que reforça o papel dos biocombustíveis como soluções de transição. No entanto, os autores ressaltam que resultados de \gls{TEA} são sensíveis a hipóteses sobre preços futuros de energia, custos de capital e evolução tecnológica, o que exige cautela na interpretação de cenários de longo prazo.

Uma limitação recorrente dos estudos baseados exclusivamente em \gls{TEA} é a ausência de avaliação ambiental integrada. Ao focar predominantemente em custos, essa abordagem pode subestimar benefícios ambientais ou externalidades negativas associadas a determinadas rotas tecnológicas. Além disso, muitos estudos assumem condições ideais de operação e escala, o que pode reduzir a aderência dos resultados, especialmente em contextos reais de aplicação, como o brasileiro. Por essa razão, a literatura recente passa a explorar abordagens que combinam desempenho econômico e ambiental em um mesmo arcabouço.

\subsection{Abordagens integradas TEA--ACV}
Com o objetivo de superar as limitações das abordagens isoladas, parte crescente da literatura adota abordagens integradas que combinam análise técnico-econômica e avaliação de ciclo de vida \citep{tan_biofuel_2021,li_techno-economic_2022,cesilla_de_souza_technical-economic_2024}. Essas abordagens permitem a avaliação simultânea de custos e impactos ambientais, possibilitando análises mais abrangentes das implicações associadas às diferentes escolhas energéticas no transporte marítimo.

\citet{tan_biofuel_2021} e \citet{li_techno-economic_2022} destacam-se como referências nesse campo ao propor estruturas de integração entre \gls{TEA} e \gls{ACV} de forma consistente, permitindo a comparação de múltiplas rotas de produção de biocombustíveis sob diferentes cenários. \citet{cesilla_de_souza_technical-economic_2024} avançam nesse sentido ao aplicar abordagem integrada para avaliar biocombustíveis navais derivados de \emph{feedstocks} brasileiros, incorporando análises de incerteza e indicadores como custo nivelado do combustível (\gls{LCOS}). Esse tipo de abordagem é relevante para o contexto brasileiro, pois permite capturar simultaneamente os efeitos ambientais e econômicos associados à disponibilidade de biomassa, à logística e à escala de produção.

Apesar de seu potencial explicativo, as abordagens integradas ainda são pouco exploradas na literatura aplicada ao transporte marítimo, especialmente quando associadas a dados operacionais reais ou simulações por rota e embarcação. A maioria dos estudos permanece no nível de cenários agregados, o que limita sua capacidade de orientar decisões operacionais específicas. Esse ponto aproxima as abordagens integradas da vertente de modelagem operacional, logística e espacial, na qual a literatura busca incorporar restrições de infraestrutura, localização e abastecimento.

\subsection{Modelagem operacional, logística e espacial}
Uma vertente metodológica menos explorada, mas de crescente relevância, envolve a modelagem operacional, logística e espacial aplicada à transição energética no transporte marítimo \citep{harahap_renewable_2023,carvalho_biofuels_2021,cesilla_de_souza_technical-economic_2024}. Em contraste com abordagens que permanecem em cenários agregados, esses estudos buscam incorporar dimensões territoriais, fluxos de abastecimento e restrições de infraestrutura, aproximando a análise de condições reais de operação.

\citet{harahap_renewable_2023} utilizam modelos de otimização, como o \emph{BeWhere}, um modelo de otimização espacial aplicado à cadeia de bioenergia, para avaliar a localização ótima de unidades de produção e distribuição de combustíveis alternativos, considerando custos logísticos e disponibilidade de matérias-primas. \citet{carvalho_biofuels_2021}, por sua vez, aplicam abordagem espacial e logística para analisar os fluxos de suprimento de biocombustíveis até potenciais pontos de produção e abastecimento portuário, evidenciando que a viabilidade do biobunker depende de fatores como distância, escala e integração com a infraestrutura portuária.

Outros estudos incorporam simulações de processo para comparar rotas de produção de combustíveis alternativos, frequentemente com foco em aspectos técnicos e de eficiência de conversão. \citet{cesilla_de_souza_technical-economic_2024}, por sua vez, integram modelagem de processo, avaliação de ciclo de vida e análise técnico-econômica em estudo aplicado ao contexto marítimo, considerando parâmetros representativos de uma embarcação. Esses trabalhos apontam para a importância de modelos integrados capazes de capturar a complexidade dos sistemas energéticos associados ao transporte marítimo, mas também evidenciam escassez de estudos que combinem modelagem operacional detalhada com dados reais de navegação.

\subsection{Comparações entre combustíveis e rotas tecnológicas}
À luz das abordagens metodológicas discutidas anteriormente, a comparação entre combustíveis fósseis, biocombustíveis e alternativas emergentes de emissão zero constitui eixo central da literatura sobre descarbonização do transporte marítimo. Essas comparações são fundamentais para avaliar não apenas o potencial de redução de emissões, mas também a viabilidade técnica, econômica e operacional das diferentes rotas tecnológicas ao longo do tempo. De modo geral, os estudos analisados podem ser organizados em três grupos de comparação: (i) combustíveis fósseis versus biocombustíveis líquidos, (ii) biocombustíveis versus combustíveis emergentes de emissão zero e (iii) comparações internas entre diferentes rotas e misturas de biocombustíveis. Em cada caso, a literatura mobiliza, em graus distintos, resultados de ensaios operacionais, avaliações de ciclo de vida, análises técnico-econômicas e estudos orientados por logística e infraestrutura.

\subsubsection{(i) Combustíveis fósseis versus biocombustíveis líquidos}
A comparação entre combustíveis fósseis convencionais e biocombustíveis líquidos é abordada na literatura, especialmente no que se refere às emissões atmosféricas e à compatibilidade operacional com motores marítimos existentes. Nesse grupo, a literatura combina evidências experimentais sobre desempenho de motores e poluentes locais com estudos de \gls{ACV} orientados à comparação de emissões ao longo do ciclo de vida. Em conjunto, esses trabalhos indicam que o uso de biodiesel e misturas com \emph{bunker fuel} resulta em reduções de poluentes locais, como óxidos de enxofre (SOx) e material particulado, além de potenciais reduções nas emissões de dióxido de carbono quando analisadas em perspectiva de ciclo de vida, a depender da matéria-prima e das condições de produção. Um dos principais fatores associados a essa variabilidade é a mudança indireta do uso da terra (\gls{ILUC}), que se refere às emissões adicionais decorrentes do deslocamento da produção agrícola para novas áreas em resposta à expansão de culturas energéticas.

\citet{petzold_operation_2011} observam reduções significativas, da ordem de até 70\%, nas emissões de material particulado e \emph{black carbon} em ensaios experimentais com motores diesel marítimos operando com combustíveis biogênicos, em comparação ao óleo combustível pesado. No entanto, os autores ressaltam que tais reduções em poluentes de curto prazo não se traduzem necessariamente em benefícios climáticos globais quando considerados os impactos do ciclo de vida do combustível, incluindo mudanças de uso da terra.

\citet{altarazi_effects_2022} reforçam esse diagnóstico ao analisar o desempenho e as emissões de motores marítimos operando com misturas de biodiesel até B20. Os autores concluem que essas misturas mantêm eficiência térmica próxima à do combustível fóssil, ao mesmo tempo em que promovem reduções expressivas nas emissões de poluentes locais, como óxidos de enxofre, monóxido de carbono, hidrocarbonetos e material particulado. Do ponto de vista operacional, esses resultados são relevantes, pois indicam que os biocombustíveis líquidos podem ser introduzidos no setor marítimo com necessidade limitada de adaptações tecnológicas, especialmente em aplicações de curto prazo.

Em perspectiva sistêmica, estudos baseados em \gls{ACV}, como os de \citet{masum_comparing_2023} e \citet{international_council_on_clean_transportation_potential_2021}, demonstram que a magnitude das reduções de emissões de \gls{GEE} associadas aos biocombustíveis depende do tipo de \emph{feedstock} e do método de produção. Rotas baseadas em resíduos tendem a apresentar melhor desempenho ambiental em comparação a combustíveis fósseis, enquanto rotas baseadas em culturas dedicadas podem apresentar resultados mais heterogêneos, especialmente quando considerados impactos indiretos relacionados ao uso da terra.

\subsubsection{(ii) Biocombustíveis versus combustíveis emergentes de emissão zero}
Outra vertente importante da literatura envolve a comparação entre biocombustíveis líquidos e combustíveis emergentes de emissão zero, como hidrogênio, amônia e \emph{e-fuels}. Essas comparações são utilizadas para discutir o papel relativo de cada rota tecnológica ao longo do horizonte temporal das metas estabelecidas pela \gls{IMO}. Aqui, resultados de \gls{TEA} e de estudos comparativos mais amplos tornam-se particularmente relevantes, pois a diferença entre as rotas não se expressa apenas em potencial máximo de abatimento, mas também em custos, maturidade tecnológica e exigências de infraestrutura.

\citet{lindstad_reduction_2021} estimam que combustíveis sintéticos e rotas baseadas em hidrogênio podem alcançar reduções de emissões superiores a 80\% em cenários de longo prazo. Contudo, os autores ressaltam que esses ganhos estão associados a custos mais elevados, à necessidade de investimentos substanciais em infraestrutura e à adoção de novos sistemas de propulsão. Além disso, limitações relacionadas ao armazenamento, à densidade energética e aos custos são apontadas como barreiras relevantes à adoção dessas alternativas no curto e médio prazo, especialmente para grandes embarcações.

Nesse contexto, \citet{kanchiralla_role_2024} argumentam que biocombustíveis e \emph{e-fuels} não devem ser tratados como rotas concorrentes, mas como soluções complementares em diferentes estágios da transição energética. Enquanto combustíveis de emissão zero representam o objetivo de longo prazo, os biocombustíveis líquidos oferecem rota de transição com potencial de redução de emissões no curto prazo, apoiada na utilização da infraestrutura e da frota existentes. \citet{mukherjee_techno-economic_2023} reforçam essa perspectiva ao analisar a competitividade dos biocombustíveis quando combinados com tecnologias de \gls{CCS} (Carbon Capture and Storage, ou captura e armazenamento de carbono), indicando que essas combinações podem ampliar o potencial de mitigação climática no setor marítimo.

\subsubsection{(iii) Comparações internas entre rotas e misturas de biocombustíveis}
Além das comparações entre grandes classes de combustíveis, diversos estudos realizam análises comparativas internas entre diferentes rotas de produção e misturas de biocombustíveis. Essas análises são relevantes para identificar quais combinações apresentam melhor desempenho ambiental e, em alguns casos, econômico. Nesse conjunto, ganham destaque comparações entre \emph{feedstocks}, rotas de conversão e métricas de avaliação, o que aproxima diretamente os resultados de debates metodológicos previamente apresentados nas subseções de \gls{ACV}, \gls{TEA} e abordagens integradas.

\citet{tan_biofuel_2021} e \citet{masum_comparing_2023} concluem que rotas baseadas em resíduos tendem a apresentar melhor desempenho ambiental, principalmente por evitarem impactos associados à expansão agrícola. Além disso, \citet{krantz_assessing_2023} mostram que a escolha da métrica de avaliação, média ou marginal, pode alterar essa percepção, influenciando a forma como o impacto climático dessas rotas é interpretado.

Estudos regionais e setoriais também contribuem para essa comparação interna. \citet{spoof-tuomi_environmental_2020} analisam a navegação de curta distância e indicam que biocombustíveis líquidos podem gerar ganhos operacionais e ambientais em contextos nos quais a infraestrutura para combustíveis emergentes é limitada. \citet{sagin_impact_2023} quantificam potenciais reduções de custos de manutenção e ganhos operacionais em motores marítimos movidos a biodiesel, enquanto \citet{doliente_bio-aviation_2020} exploram interseções logísticas entre os setores marítimo e aéreo, sugerindo que sinergias na cadeia de suprimento podem influenciar a viabilidade de determinadas rotas de biocombustíveis.

\subsection{Síntese das evidências e identificação de lacunas de pesquisa}
A análise dos trabalhos relacionados evidencia que a literatura sobre descarbonização do transporte marítimo avançou significativamente na identificação e avaliação de rotas tecnológicas alternativas ao \emph{bunker fuel}. As evidências revisadas convergem ao reconhecer que os biocombustíveis líquidos apresentam potencial relevante para a redução de emissões e que sua atratividade depende do \emph{feedstock}, da rota de produção, das escolhas metodológicas e das condições logísticas e institucionais de aplicação. Do ponto de vista metodológico, observa-se ampla utilização de abordagens de \gls{ACV} e \gls{TEA}, com tendência recente de integração entre essas vertentes, embora ainda predomine o uso de análises agregadas, com limitada incorporação de dados operacionais reais e restrições específicas de rotas, embarcações e infraestrutura portuária.

Apesar desses avanços, a revisão evidencia lacunas importantes. Primeiramente, há escassez de estudos que integrem, de forma consistente, avaliação ambiental, análise econômica e modelagem operacional aplicada ao transporte marítimo. Poucos trabalhos combinam \gls{ACV}, \gls{TEA} e simulação logística em um mesmo arcabouço, o que limita a capacidade de avaliar \emph{trade-offs} entre custos, emissões e viabilidade operacional. Em segundo lugar, observa-se concentração geográfica dos estudos em contextos europeus e globais, com limitada atenção às especificidades de países em desenvolvimento, particularmente ao caso brasileiro, no qual a disponibilidade de biomassa, a estrutura portuária e os instrumentos de política pública diferem substancialmente.

No caso brasileiro, a literatura é incipiente no que se refere à aplicação de modelos integrados que utilizem dados reais de operação marítima e considerem explicitamente a logística de abastecimento, a localização de fontes de biomassa e o arcabouço institucional vigente \citep{carvalho_biofuels_2021,grangeia_brazilian_2022,cesilla_de_souza_technical-economic_2024}. Embora existam estudos relevantes sobre biocombustíveis e políticas públicas, como o RenovaBio, poucos trabalhos exploram sua aplicação direta ao transporte marítimo e ao uso de biocombustíveis como biobunker \citep{brasil_lei_2017,grangeia_brazilian_2022}. Ademais, não foram identificados estudos nacionais que desenvolvam simuladores operacionais capazes de otimizar misturas de biocombustíveis por rota ou embarcação, considerando simultaneamente custos, emissões e condicionantes institucionais e regulatórios.

Essas lacunas fundamentam a proposta desta dissertação, que investiga o biobunker como combustível de transição no transporte marítimo internacional operando a partir do Brasil por meio de abordagem integrada. O estudo articula dados operacionais reais, avaliação ambiental baseada em ciclo de vida e análise econômica, considerando as restrições técnicas, logísticas e institucionais associadas a essa rota energética. Ao adotar esse enquadramento, a dissertação busca contribuir para a compreensão das possibilidades e dos limites do biobunker no contexto brasileiro, especialmente no horizonte das metas intermediárias de descarbonização estabelecidas pela \gls{IMO} para 2030, sem assumir caráter normativo ou prescritivo.
